
Pretrained video generation models~\cite{dreamzero,wan2025wan,hu2024vpp,liang2025videogenerators} have emerged as a promising foundation for robot learning because they capture rich visual dynamics, object interactions, and multi-step temporal structure from large-scale video data. In parallel, Vision-Language-Action (VLA) models have shown strong progress in end-to-end robot control by directly mapping visual observations and language instructions to actions~\cite{rt2,palme,openvla,atomvla,llava-vla,pi0,pi05}. Motivated by the strengths of both lines, recent unified world action models seek to jointly predict future observations and future actions within a single generative framework, offering a path toward general-purpose robot control that combines visual foresight with action generation~\cite{unified_wm,li2026causal,bi2025motusunifiedlatentaction,pertsch2025fast,gigaworld,dreamzero}.

Despite this promise, decoding reliable actions from such models still requires substantial robot-domain adaptation. We argue that a key limitation is structural rather than purely statistical. Predicting future RGB frames answers what the scene may look like, but control additionally depends on where the robot should interact and why that interaction is useful for the task. This type of spatially grounded reasoning has long been recognized as important in robotic manipulation~\cite{shridhar2022cliport,mo2021where2act,shridhar2023peract,wi2023calamari}. In cluttered manipulation, these signals are sparse, while future RGB representations are dense and dominated by appearance details that are often irrelevant to the next motor command. As a result, action decoding directly from future visual latents forces the model to recover manipulation intent implicitly from a representation that is not optimized for control.

To bridge this gap, we introduce \textbf{\method}, an intent-aware unified world action model that establishes an explicit spatial interface between future prediction and action generation. Rather than decoding actions directly from future visual representations, \method predicts an aligned action-based spatial value map that encodes the interaction structure required for task execution. This intermediate representation transforms future scene dynamics into a control-oriented spatial prior, providing a more direct pathway for action generation. Built on a shared pretrained video prior, \method jointly models future frames and future value maps, while an \emph{intent-causal attention} mechanism forces the action branch to access future information only through the predicted value map rather than directly through future visual tokens. In this way, the model separates scene evolution from manipulation intent without sacrificing their alignment.

We further develop a self-distillation RL post-training stage that freezes the video and value branches and optimizes only the action head using dense rewards induced by projected value-map responses together with sparse task-level rewards. To support training and evaluation, we build a large-scale simulation dataset with synchronized multi-view videos, action sequences, and value-map annotations. Experiments on the RoboTwin 2.0 benchmark~\cite{robotwin2.0} show that explicit spatial intent modeling consistently improves unified world action learning over strong baselines, while also yielding a more interpretable action-generation process through localized stage-wise interaction regions.

In summary, this work makes three main contributions:
\begin{itemize}[leftmargin=*]
    \item \textbf{Intent-aware unified world action model.}
    An explicit spatial value-map interface is introduced between future prediction and action decoding, enabling the model to capture task-relevant interaction structure in a control-oriented form.

    \item \textbf{A unified training framework for spatially grounded control.}
    Joint frame-value generation, intent-causal attention, and post-training self-distillation RL are integrated into a single framework that improves action learning without disturbing the prior training video.

    \item \textbf{Large-scale simulation dataset and strong benchmark performance.}
    We construct a large-scale simulation dataset of 30K manipulation trajectories with synchronized multi-view videos, actions, and value-map annotations. On the RoboTwin 2.0 benchmark~\cite{robotwin2.0}, AIM achieves \textbf{94.0\%} and \textbf{92.1\%} average success rates under the Easy and Hard settings, respectively, while consistently outperforming prior external baselines.
\end{itemize}


